\documentclass[10pt,a4paper]{amsart}
\usepackage[margin=19mm]{geometry}
\usepackage{amsmath,amssymb,mathtools}
\usepackage[scr=rsfs,cal=euler]{mathalfa}
\usepackage{xfrac}
\usepackage[unicode,hidelinks,bookmarks=false]{hyperref}
\newcommand{\ie}{i.e.\ }
\newcommand{\cf}{cf.\ }
\newcommand{\resp}{respectively\ }
\newcommand{\Subsection}{\subsection}
\providecommand{\nobreakdash}{\nobreak\hbox{-}}
\setlength{\emergencystretch}{2em}
\multlinegap0pt
\usepackage{mathtools}
\usepackage{bm}
\usepackage{float}
\usepackage{xcolor}
\usepackage[normalem]{ulem}
\usepackage[shortlabels]{enumitem}
\usepackage{tcolorbox}
\newtheorem{lemma}{Lemma}
\newtheorem{pf}{Pf}
\newtheorem{fleqn}{Fleqn}
\usepackage{cleveref}
\crefname{lemma}{Lemma}{Lemmas}
\crefname{pf}{Pf}{Pfs}
\crefname{fleqn}{Fleqn}{Fleqns}
\title{Examining the impact of forcing function inputs on structural identifiability}
\author{Jessica R Conrad, James M Hyman, Marisa C Eisenberg}
\date{Source: arXiv:2407.02771v2 (CC BY 4.0)}
\begin{document}
\maketitle

\noindent\emph{Source and licence.} Examining the impact of forcing function inputs on structural identifiability. Version arXiv:2407.02771v2 (CC BY 4.0), distributed by arXiv under the Creative Commons Attribution 4.0 International licence, http://creativecommons.org/licenses/by/4.0/, as recorded in the arXiv licence metadata for this exact version and shown on the abstract page https://arxiv.org/abs/2407.02771v2. Full licence notice: \texttt{licenses/arxiv-2407.02771-CC-BY-4.0.txt}.\par\smallskip

\noindent\textit{Editorial scope.} Counted unit: the article's lemma infsafe (source0.tex lines 679--684). The article's complete original proof of it is delivered with it (source0.tex lines 686--750, one \texttt{proof} environment of 65 lines). No mathematics is written, repaired or re-derived: the statement and the proof are the article's own lines, in the article's own notation, and no retained line is substituted or re-flowed.  Retained uncounted as the source-local context the proof needs: lines 271--277; lines 492--494; lines 620--631.  Nothing was omitted from any retained span, so the package carries no omission record.  No retained line contains a citation, so the wrapper carries no bibliography and no bibliography entry.  No retained line contains a figure environment, an image-inclusion command or drawing code, so no image is delivered and the package declares no asset dependency.  Editorial formatting only, in addition: an AMS \texttt{amsart} preamble that restates the article's own declarations and macros used by the retained lines, namely the environments \texttt{lemma}, \texttt{pf}, \texttt{fleqn} on separate counters, together with the standard packages those lines need.  The retained environments print in this document as Lemma 1, Pf 1, Fleqn 1 in order of appearance, whereas the article prints the same items with the numbering of its own full document.  The complete native source of the article as distributed in the arXiv e-print for this exact version is bundled unchanged as \texttt{sources/161743/source0.tex}.
\begin{equation} \label{eq:model}
    \begin{aligned}
    \dot{\pmb{x}} &=f(\pmb{x}(t, \pmb{\theta}),\pmb{u}(t),\pmb{\theta}),\\
    \pmb{x}(0, \pmb{\theta}) &= x_0(\pmb{\theta}), \\
    \pmb{y}(t, \pmb{\theta}) &= g(\pmb{x}(t, \pmb{\theta}), \pmb{\theta}),
    \end{aligned}
\end{equation}

\begin{multline} \label{ranking}
u_1 < u_2 < ... <\dot{u}_1 < \dot{u}_2 < ... \\< y_1 < y_2 < ... <\dot{y}_1 < \dot{y}_2 < ...\\< x_1 < x_2 < ...<\dot{x}_1 < \dot{x}_2 < ...
\end{multline}

\begin{lemma} \label{newform} 
Let $\tilde{u}$ be a forcing input function, and let $M$ be any rational function model of the form of \cref{eq:model}, assuming the standard ranking described in \cref{ranking}. Generate $\tilde{M}$ by scaling any parameter $\tilde{\theta} \in \pmb{\theta}$ in the state variable equations of $M$ with $\tilde{u}$.
     
Then for each input-output relation precursor equation $P'(M)$, $\tilde{M}$ will have an analogous input-output equation precursor $P'(\tilde{M})$, where the terms of $P'(M)$ will also be present in $P'(\tilde{M})$ with two modifications: 1) the replacement rule $\tilde{\theta} \mapsto \tilde{\theta}\tilde{u}$ (which may split a monomial into two or more distinct terms), and 2) the multiplication of all such analogous terms in $P'(\tilde{M})$ by a common term in the parameters and $\tilde{u}$. We describe these modifications in more detail below.

Let $C_im_i = (c_0 + c_1\tilde{\theta} + \cdots c_n\tilde{\theta}^n)m_i$ be a monomial term in $P'(M)$, where $C_i$ is a coefficient polynomial in the parameters (potentially including $\tilde{\theta}$, written here as a polynomial in $\tilde{\theta}$ where some $c_i$ may be zero), and  $m_i$ is a monomial in outputs, inputs, and their derivatives. 

Then $P'(\tilde{M})$ will contain analogous terms to $C_im_i$, of the form $ac_0m_i$, $ac_1\tilde{\theta}\tilde{u}m_i$, $\dots$ $ac_n\tilde{\theta}^n\tilde{u}^nm_i$, where $a$ is a common multiplicative term given by a polynomial in the parameters multiplied by a monomial in $\tilde{u}$ ($a$ may be trivial, e.g. $a=1$).

Additional terms not descended from any original term may also be generated in $P'(\tilde{M})$, but these will never be the leading term of $P'(\tilde{M})$.

\end{lemma}

\begin{lemma} \label{infsafe}
    Let $P'(M)$ be the precursor to an input-output equation $P(M)$ derived from a rational function model $M$ of the form \cref{eq:model}, assuming the standard ranking described in \cref{ranking}. 
    Let $\tilde{M}$ be the scaled model following the replacement $\tilde{\theta}\mapsto \tilde{\theta}\tilde{u}$ for some $\tilde{\theta} \in \pmb{\theta}$, and let $P(\tilde{M})$ be the resulting corresponding input-output equation to $P(M)$ (i.e. such that $P'(\tilde{M})$ contains the corresponding monomials given in Lemma 1).
    
    Then the parameter solutions for $P(\tilde{M})$ are contained within the parameter solutions for $P(M)$. 
\end{lemma}

\begin{pf}
    Consider $P'(M)$, associated with some input-output equation $P(M)$ derived from model $M$. 
    $P'(M)$ has the general form $0 = P'(M) = \sum_{i = 1}^n C_i m_i$, where the $C_i$ are polynomials in the parameters, the $m_i$ are monomials in the inputs, outputs, and their derivatives, and $m_n$ is the \textit{leading monomial} (noting also that $n$ here is just a positive integer, not the number of state variables in the system).
    
    Each coefficient $C_i$ can itself be viewed as a polynomial in $\tilde{\theta}$ with coefficients in $\pmb{\theta}\setminus\tilde{\theta}$. We will write this as $C_i = c_i + \tilde{c}_i \tilde{\theta}^{p_i}$, where  $\tilde{c}_i \tilde{\theta}^{p_i}$ is the highest degree term in $\tilde{\theta}$ (possibly with $\tilde{c}_i=0$ if there is no $\tilde{\theta}$ present in that $C_i$), and $c_i$ is all remaining terms in $C_i$. $P'(M)$ is then: $0 = P'(M) = \sum_{i = 1}^n (c_i + \tilde{c}_i \tilde{\theta}^{p_i}) m_i = (c_1 + \tilde{c}_1 \tilde{\theta}^{p_1}) m_1 + ... + (c_{n-1} + \tilde{c}_{n-1} \tilde{\theta}^{p_{n-1}}) m_{n-1} + (c_{n} + \tilde{c}_{n} \tilde{\theta}^{p_{n}}) m_{n}$.

    Next, we make $P'(M)$ a monic polynomial $P(M)$ by normalizing with respect to the leading coefficient of $m_n$:
    \begin{fleqn}
    \begin{align*}
        0 = P(M) = &\frac{c_1 + \tilde{c}_1 \tilde{\theta}^{p_1}}{c_{n} + \tilde{c}_{n} \tilde{\theta}^{p_{n}}} m_1 + ... \\&~~~~~+ \frac{c_{n-1} + \tilde{c}_{n-1} \tilde{\theta}^{p_{n-1}}}{c_{n} + \tilde{c}_{n} \tilde{\theta}^{p_{n}}} \tilde{m}_{n-1} + m_n
    \end{align*}
    \end{fleqn}
    which now satisfies the definition of an input-output equation. The parameter solutions of $P(M)$ must therefore satisfy:
    \begin{align}
        \frac{c_i + \tilde{c}_i \tilde{\theta}^{p_i}}{c_n + \tilde{c}_n \tilde{\theta}^{p_n}} & = \frac{c^*_i + \tilde{c}_i^* \tilde{\theta}^{*~p_i}}{c_n^* + \tilde{c}_n^* \tilde{\theta}^{*~p_n}}, ~~~\text{for}~i=1,...,n-1 \label{psoln_thm1}
    \end{align}
    where $c_j^*$ and $\tilde{\theta}^*$ represent the true parameter values.
    
    Similarly, $P'(\tilde{M})$ is a polynomial of at least $n$ monomials now calculated with $\tilde{\theta}$ scaled by a known forcing input function $\tilde{u}$. 
    By Lemma \ref{newform}, $P'(\tilde{M})$ has the general form: 
     \begin{align*}
        0 = P'(\tilde{M}) & =  \sum_{i = 1}^{n} a(c_i + \tilde{c}_i \tilde{\theta}^{p_i}\tilde{u}^{p_i} )  \tilde{m}_i + \sum_{j = 1}^{k} a_j \\
        & =  a(c_1 + \tilde{c}_1 \tilde{\theta}^{p_1} \tilde{u}^{p_1}) \tilde{m}_1 + \cdots \\
        & ~~~~~~~~+ a(c_{n} + \tilde{c}_{n} \tilde{\theta}^{p_{n}}\tilde{u}^{p_{n}}) \tilde{m}_{n} + \cdots
    \end{align*}
    where $\tilde{m}_i$ contain $m_i$ along with any additional monomial of $\tilde{u}$, and $a$ is a multiplicative coefficient in the parameters. The $a_j$ are any additional terms also generated (which by Lemma \ref{newform} are not the leading term).
    
    To make $P'(\tilde{M})$ a monic polynomial $P(\tilde{M})$, we normalize with respect to the coefficient of the leading monomial in $P'(\tilde{M})$. 
    
    If $\tilde{c}_n \neq 0$, then the leading monomial of $P(\tilde{M})$ is $\tilde{u}^{p_n}m_n$, and its coefficient is $a\tilde{c}_n \tilde{\theta}^{p_n}$.
    Then the terms in $P(\tilde{M})$ corresponding to the original monomials $m_i$ (now possibly split) are of the form $\frac{c_i}{\tilde{c}_n \tilde{\theta}^{p_n}} \tilde{m}_i$ and $\frac{\tilde{c}_i \tilde{\theta}^{p_i}}{\tilde{c}_n \tilde{\theta}^{p_n}} \tilde{u}^{p_i} \tilde{m}_i$.
    The parameter solutions of $P(\tilde{M})$ must therefore satisfy:
    \begin{align}
        \frac{c_i}{\tilde{c}_n \tilde{\theta}^{p_n}} & = \frac{c^*_i}{\tilde{c}_n^* \tilde{\theta}^{*~p_n}}, ~~~\text{for}~i=1,...,n \label{pnewsoln_thm1} \\
        \frac{\tilde{c}_i \tilde{\theta}^{p_i}}{\tilde{c}_n \tilde{\theta}^{p_n}} & = \frac{\tilde{c}^*_i \tilde{\theta}^{*~p_i}}{\tilde{c}_n^* \tilde{\theta}^{*~p_n}}, ~~~\text{for}~i=1,...,n \label{pnewsolnB_thm1} 
    \end{align} 
   
    We want to show that the parameter solutions of $P(\tilde{M})$ are also parameter solutions for $P(M).$
    From the parameter solutions \cref{pnewsoln_thm1} and \cref{pnewsolnB_thm1} of $P(\tilde{M})$, we have:
    \begin{align*}
        c_i & = c^*_i \frac{\tilde{c}_n \tilde{\theta}^{p_n}}{\tilde{c}_n^* \tilde{\theta}^{*~p_n}}, & ~~\text{for~}i=1,...,n \\
        \tilde{c}_i \tilde{\theta}^{p_i} & = \tilde{c}_i^* \tilde{\theta}^{*~p_i} \frac{\tilde{c}_n \tilde{\theta}^{p_n}}{\tilde{c}_n^* \tilde{\theta}^{*~p_n}}, & ~~\text{for~}i=1,...,n
    \end{align*}
    Substituting these into the left-hand side of \cref{psoln_thm1}:
    \begin{fleqn}
    \begin{align*}
        \frac{c_i + \tilde{c}_i \tilde{\theta}^{p_i}}{c_n + \tilde{c}_n \tilde{\theta}^{p_n}} &= \frac{c^*_i \frac{\tilde{c}_n \tilde{\theta}^{p_n}}{\tilde{c}_n^* \tilde{\theta}^{*~p_n}} + \tilde{c}_i^* \tilde{\theta}^{*~p_i} \frac{\tilde{c}_n \tilde{\theta}^{p_n}}{\tilde{c}_n^* \tilde{\theta}^{*~p_n}}}{c_n^* \frac{\tilde{c}_n \tilde{\theta}^{p_n}}{\tilde{c}_n^* \tilde{\theta}^{*~p_n}} + \tilde{c}_n^* \tilde{\theta}^{*~p_n} \frac{\tilde{c}_n \tilde{\theta}^{p_n}}{\tilde{c}_n^* \tilde{\theta}^{*~p_n}}} \\
        &= \frac{(c^*_i + \tilde{c}_i^* \tilde{\theta}^{*~p_i}) \frac{\tilde{c}_n \tilde{\theta}^{p_n}}{\tilde{c}_n^* \tilde{\theta}^{*~p_n}}}{(c_n^* + \tilde{c}_n^* \tilde{\theta}^{*~p_n}) \frac{\tilde{c}_n \tilde{\theta}^{p_n}}{\tilde{c}_n^* \tilde{\theta}^{*~p_n}}} = \frac{c^*_i + \tilde{c}_i^* \tilde{\theta}^{*~p_i}}{c_n^* + \tilde{c}_n^* \tilde{\theta}^{*~p_n}}
    \end{align*}
    \end{fleqn}
    This matches the right-hand side of \cref{psoln_thm1}. Therefore the solutions are self-consistent.
    We conclude that solutions of $P(\tilde{M})$ are also parameter solutions for $P(M)$. Note that this calculation becomes slightly more complicated if the $c_i$ include lower powers of $\tilde{\theta}$ (as then the $c_i$ split into distinct monomials), but the basic argument still holds. 
    
    In the case when $\tilde{c}_n = 0$, the coefficient $C_n$ does not contain $\tilde{\theta}$, so after the replacement $\tilde{\theta} \mapsto \tilde{\theta}\tilde{u}$, we have $C_n$ unchanged. The leading term coefficient in both $P(M)$ and $P(\tilde{M})$ is therefore $c_n$ (the highest-degree term in $C_n$ when $\tilde{c}_n = 0$).

After normalization, the parameter solutions must satisfy:
\begin{align*}
\frac{c_i}{c_n} &= \frac{c^*_i}{c_n^*}, \quad \text{and} \quad \frac{\tilde{c}_i \tilde{\theta}^{p_i}}{c_n} = \frac{\tilde{c}_i^* \tilde{\theta}^{*~p_i}}{c_n^*}
\end{align*}
for $\tilde{M}$, and:
\begin{align*}
\frac{c_i + \tilde{c}_i \tilde{\theta}^{p_i}}{c_n + \tilde{c}_n \tilde{\theta}^{p_n}} = \frac{c_i + \tilde{c}_i \tilde{\theta}^{p_i}}{c_n} = \frac{c^*_i + \tilde{c}_i^* \tilde{\theta}^{*~p_i}}{c_n^*}
\end{align*}
for $M$ (since $\tilde{c}_n = 0$). The containment of solution sets follows immediately by adding the two equations from $\tilde{M}$. $\square$ 
\end{pf}

\end{document}
